% Page 009 — page imprimée 5
% Trois figures marquantes à Meyrueis
% Reconstruction éditable en LaTeX.
%
% La page se termine au milieu d'une phrase : la suite est sur la page suivante.


\begin{center}
  {\Large\bfseries Trois figures marquantes à Meyrueis}
\end{center}

\vspace{0.65em}

Grâce à d’assez nombreuses archives et études déjà menées, ( en particulier par Pierre Rolland Charles Bost, ainsi que Philippe Chambon et Patrick Cabanel ) j’ai pu retracer la vie et les activités de trois hommes ; Jacques Crégut de Ferrussac, le camisard, Jean Roman le prédicant, et Pierre Bertrand de Campis, le galérien.

\vspace{0.3em}

{\large\bfseries Jacques Crégut, le camisard}

C’est à Ferrussac qu’il est né vers 1678. Sa famille avait été gagnée à la Réforme comme toutes les familles des vallées de Meyrueis. Paysan, tisserand de cadis il allait se louer dans les basses vallées des Cévennes pour les travaux des champs ou le ramassage des châtaignes. C’est au cours d’une de ses « louées » qu’il fut enrôlé de force dans le régiment en garnison dans la région, par un agent recruteur. Il déserta et rencontra la troupe de Camisards dirigée par Castanet de Massevaques. Le connaissait-il ? en avait-il entendu parler ? c’est probable. Il le suivit et devint camisard. Arrêté du côté d’Aulas, en 1703, il sauva sa vie en donnant tous les renseignements demandés, y compris le nom de certains de ses camarades, à M. de Ginestous, subdélégué de Mgr de Basville, Intendant du Languedoc.

\begin{center}
  {\large Interrogation de Jacques Crégut}

  \vspace{0.55em}

  \textbf{Par Henry de Ginestous, subdélégué de Monseigneur de Basville, Intendant du Languedoc\\
  Le 23 Avril 1703, prison du Vigan}
\end{center}

\vspace{0.35em}

« Nous avons mandé venir des prisonniers, un jeune homme de moyenne taille, ayant les cheveux noirs et assez longs, étant en chemise de cadis-burel, portant deux paires de chausses l’une gris l’autre burel, portant une chemise de toile blanche, portant des bas faits à l’aiguille, blancs, et de bons souliers et un chapeau retroussé. »

\begin{quote}
\itshape
auquel nous avons demandé son nom, son surnom, sa demeure, sa religion, son âge, sa qualité et le sujet de sa détention :
\end{quote}

Nous a dit se nommer Jacques Crégut, tisserand de cadis fils de feu Jean du massage \textit{(hameau)} de Ferrussac, âgé d’environ 25 ans, faisant profession de la Religion Prétendue Réformée. Huit jours après la Saint Michel il part à Aulas s’étant loué chez Alibert, hôtelier, de la treille du Goulard ( ? ) pour amasser les châtaignes, avec François Martin, fournier du Goulard, il va à une assemblée qui se fait au dessus de Salagosse dans une claye \textit{(clède : séchoir à châtaignes)} (appelée la Borie du Gendre) et ayant environ 80 personnes, hommes et femmes qu’il ne connaît pas.. Le sieur François lui ayant dit que le prédicant s’appelait Delenne est d’une taille un peu plus grande que lui, avec les cheveux châtains.

Huit jours avant la fête de Noël dernier il aurait quitté le Goulard étant allé travailler de son métier de tisserand de cadis jusqu’à la fin du mois de janvier dernier que le fils du Sieur

